\section{Proofs for admissible integer classes}
\label{app:lattice}

\subsection{Hemispace separation and semantic attainment}
\label{app:lattice:semantic}

We first prove the separation fact needed when a projected convex value
function has a nonopen strict sublevel set. A hemispace in an affine space
$V$ is a convex subset of $V$ whose complement in $V$ is also convex. It
may assign different parts of a supporting hyperplane to its two sides.

\begin{lemma}[Hemispace separation]
\label{lattice:hemispace}
If $A,B$ are disjoint convex subsets of a finite-dimensional affine space
$V$, there is a hemispace $G\subseteq V$ with $A\subseteq G$ and
$G\cap B=\varnothing$.
\end{lemma}

\begin{proof}
We prove the assertion by induction on $\dim V$. If $A=\varnothing$,
take $G=\varnothing$; if $B=\varnothing$, take $G=V$. These cases also
settle dimension zero. Suppose that both are nonempty and $\dim V>0$.

Identify the direction space of $V$ with a Euclidean space. There is a
nonzero linear functional $g$ on it and a finite $\alpha$ such that
$g(a)\ge\alpha\ge g(b)$ for all $a\in A,b\in B$. Here is the
finite-dimensional separation argument, including its boundary case.
The convex set $C=A-B$ does not contain zero. If
$0\notin\cl C$, the nearest point $p$ of the nonempty closed convex set
$\cl C$ to zero satisfies
$\langle p,c\rangle\ge\|p\|^2>0$ for every $c\in C$.
If $0\in\cl C$ and $\cl C$ lies in a proper linear subspace, any nonzero
functional orthogonal to that subspace gives $g(c)=0$.
In the remaining case $\cl C$ is full-dimensional, and zero is a boundary
point of $\cl C$. To justify the latter assertion, an interior point of
the closure of a full-dimensional convex set lies in that set: choose a
small simplex around the point in the interior of the closure, approximate
its vertices by points of the convex set, and keep the point in the
interior of the approximating simplex. Convexity then includes the point.
Since $0\notin C$, it cannot be interior to $\cl C$.
Choose $q_j\notin\cl C$ tending to zero, and let $p_j$ be its nearest
point in $\cl C$. Since $0\in\cl C$, $p_j\to0$. The unit vectors
$(p_j-q_j)/\|p_j-q_j\|$ have a convergent subsequence, with unit limit
$v$. The projection inequalities
$\langle p_j-q_j,c-p_j\rangle\ge0$ yield
$\langle v,c\rangle\ge0$ for all $c\in C$. Thus in every case there is
a nonzero $g$ with $g(a)\ge g(b)$ for all $a,b$. Fixing one point of each
set shows that $\sup_B g$ and $\inf_A g$ are finite; choose $\alpha$
between them. On the affine space, $g$ is interpreted as a nonconstant
affine functional after choosing an origin.

Let $H=\{g=\alpha\}$. It has dimension one less than $V$. By induction,
there is a hemispace $G'$ of $H$ containing $A\cap H$ and missing
$B\cap H$. Put
\begin{equation*}
 G=\{g>\alpha\}\cup G'.
\end{equation*}
For a strict convex combination of two points of $G$, either one point
has $g>\alpha$ and the combination does too, or both lie in the convex set
$G'$. Hence $G$ is convex. Its complement is
$\{g<\alpha\}\cup(H\setminus G')$, which is convex by the same argument.
The points of $A$ either have $g>\alpha$ or belong to $A\cap H\subseteq
G'$. The points of $B$ have $g\le\alpha$, and those on $H$ miss $G'$.
This proves the induction.
\end{proof}

\begin{proof}[Attainment in \cref{lattice:class-number}]
Take an admissible partition $I_1,\ldots,I_k$. The strict sublevel set
$E_\tau=\{x:\phi(x)<\tau\}$ is convex; for $\tau=+\infty$ it is the
effective domain of the proper projected function. Admissibility means
$\operatorname{conv}I_j\cap E_\tau=\varnothing$. A depth-one tree whose
children are these convex hulls covers $P$ and has $k$ certified leaves.

For a binary tree, use \cref{lattice:hemispace} to find hemispaces
$G_j\supseteq\operatorname{conv}I_j$ with $G_j\cap E_\tau=\varnothing$.
When $k=1$, $G_1$ itself is a valid root and a certified leaf. It contains
$P=I_1$ and gives one node.

When $k\ge2$, begin with any convex root $N_0$ containing $P$. For
$j=1,\ldots,k-2$, split $N_{j-1}$ into
\begin{equation*}
 N_{j-1}\cap G_j\quad\text{and}\quad
 N_j=N_{j-1}\cap(\R^n\setminus G_j).
\end{equation*}
The first child is a leaf and the second continues the chain. Both are
convex, nested in their parent, and together cover that parent. At the
last internal node $N_{k-2}$, use the two leaves
$N_{k-2}\cap G_{k-1}$ and $N_{k-2}\cap G_k$. To check their required
coverage, take $p\in P\cap N_{k-2}$. It is outside $G_1,\ldots,G_{k-2}$,
and hence outside $I_1,\ldots,I_{k-2}$. Because the $I_j$ partition $P$,
$p\in I_{k-1}\cup I_k\subseteq G_{k-1}\cup G_k$. This verifies
\eqref{lattice:integer-cover}, including its restriction to integer
points at the final split. Every leaf lies in a $G_j$, so it misses
$E_\tau$ and has bound at least $\tau$. The full binary chain has $k$
leaves and $2k-1$ nodes. If a constructed leaf were empty of required
points, routing to the others would contradict minimality of $k$.

For finite $P$, replace each node set $S$ by
$S'=\operatorname{conv}(S\cap P)$. This is a polytope or the empty set.
Convexity gives $S'\subseteq S$, while
$S'\cap P=S\cap P$: one inclusion follows from $S'\subseteq S$, the
other because every point of $S\cap P$ is among the hull's generators.
Thus every node's required points and every coverage relation are
unchanged. Effective sets only shrink, so leaf bounds cannot decrease.
This gives the claimed polyhedral certificate.
\end{proof}

If $E_\tau$ is nonempty and open, its disjoint convex class hulls can be
separated by closed halfspaces. Indeed, weak separation gives
$g^Tx\ge\alpha$ on the class hull and $g^Tx\le\alpha$ on $E_\tau$;
openness and $g\ne0$ make the latter inequality strict. Thus a closed
halfspace containing the hull misses $E_\tau$. If $\phi$ is differentiable
and its minimum on a class hull is attained at $\widehat x$, first-order
optimality gives the explicit halfspace
\begin{equation*}
 \{x:\nabla\phi(\widehat x)^T(x-\widehat x)\ge0\}.
\end{equation*}
Convexity puts its value at least $\phi(\widehat x)\ge\tau$ throughout
that halfspace. If the gradient is zero, the entire space has value at
least $\tau$, so the class number is one. These special forms do not
replace the hemispace proof when $E_\tau$ is nonopen.

\subsection{Operation accounting}
\label{app:lattice:operations}

\begin{proof}[Proof of \cref{lattice:operations}]
Normalize the run's node sets to their effective sets. For part (a), route
required integer points to leaves as in the class-number lower-bound
proof. A routed point is required at every ancestor. It belongs to each
convex cut added on its path, by cut validity at the node where it was
added. Thus the convex hull of the points assigned to a leaf is contained
in the leaf's relaxation domain after all such cuts. The pruning bound
makes that hull admissible for the original $\phi$. The same argument
applies to a convex reduction preserving the required integer points. If
the cuts preserve only feasible integer projections, use $P=F$; it is not
legitimate to retain a larger required set without its preservation
hypothesis.

For part (b), replace each reduction sequence at a node by a chain. At
stage $i$, let $Q^{(i-1)}$ have the two children
$Q^{(i-1)}\cap D_i$ and $Q^{(i)}$. The first is a certified leaf, and the
second continues. Both are convex and nested, and
\eqref{lattice:removal-hypothesis} is precisely their integer-coverage
condition. Attach the run's original children to the last retained set,
or make it a leaf if the run pruned it. The expanded tree has at most
$L+S$ leaves and is a $P$-covering $\tau$ certificate. Applying
\cref{lattice:class-number} proves \eqref{lattice:augmented-count}.
Simultaneous certificates instead give extra leaf children at the
unreduced node, with the same count.

In one pass, at most $2n$ bounds change per node, so
$\kappa_\tau(P)\le L+2nN\le(2n+1)N$. A pass at the root alone gives
$\kappa_\tau(P)\le L+2n$.

For binaries, consider integer-rounded changes which remove a remaining
binary value. Before a node empties, each free variable can be fixed only
once, giving at most $n$ removals. A further change of an already fixed
variable empties the node; such a run node is a leaf and has at most
$n+1$ removals. Its final retained set has no required integer points, so
the expansion can drop that retained leaf: at the last step all required
points belong to the certified removed child. If $e$ run leaves are
emptied, the expanded leaf count is at most
$L+S-e\le L+nN$. Hence
$\kappa_\tau(P)\le L+nN\le(n+1)N$.

For general integer bounds, each genuine rounded change shortens the sum
of the integer widths by at least one while the node stays nonempty.
There can be up to $\sum_i(u_i-\ell_i)$ such changes and one final
emptying change at that node. This depends on its widths rather than only
its dimension, and supplies no dimension-only estimate for iterated
passes.
\end{proof}

If an enumeration node creates children only for a finite integer range
$a\le\pi^Tx\le b$, and its omitted values have relaxation bound at least
$\tau$, append the two certified tails
$\pi^Tx\le a-1$ and $\pi^Tx\ge b+1$. This converts the run into a
$P$-covering certificate. There are at most two added leaves per internal
node, so its expanded leaf count is at most $L+2N_{\rm int}\le3N$.
Consequently the original run has at least $\kappa_\tau(P)/3$ nodes.
The validity of the tail bounds is an assumption about the enumeration's
radius tests, not an automatic consequence of skipping a value.

\subsection{The lattice-free facet interpretation}
\label{app:lattice:facets}

\begin{proof}[Proof of \cref{lattice:facet-number}]
A finite convex function on $\R^n$ is continuous. Thus the nonempty strict
sublevel set $E_\tau$ is open and convex. Each admissible class hull is
disjoint from it. The open-set separation argument above supplies a
closed halfspace containing that hull and missing $E_\tau$. An admissible
partition therefore gives a halfspace cover of the integer lattice.
Conversely, the integer points of any closed halfspace missing $E_\tau$
are admissible, because their hull stays in the halfspace. This identifies
the first two minima.

Write a halfspace cover as
$H_j=\{x:g_j^Tx\ge\alpha_j\}$, with $g_j\ne0$. Its open complementary
polyhedron
\begin{equation*}
 O=\{x:g_j^Tx<\alpha_j\text{ for all }j\}
\end{equation*}
contains $E_\tau$ and has no integer point. Because it is nonempty and
open, its closure is
$M=\{x:g_j^Tx\le\alpha_j\text{ for all }j\}$ and
$\operatorname{int}M=O$. To see closure equality, fix $x_0\in O$; the
segment from any point of the closed system to $x_0$ satisfies every row
strictly except possibly at its first endpoint. Each genuine facet of
$M$ comes from one of these rows, so there are no more facets than
covering halfspaces. Conversely, the closed complementary halfspaces of
the facets of any full-dimensional polyhedron satisfying
\eqref{lattice:enclosing-polyhedron} cover $\Z^n$ and miss $E_\tau$.
The minima therefore agree.

For the bound $2^n$, use the maximal lattice-free polyhedron theorem
\citep[Theorem~2]{basuConfortiCornuejolsZambelli2010MaximalLatticeFree}:
every nonempty open convex set avoiding $\Z^n$ is contained in a
full-dimensional maximal lattice-free convex set; every such maximal set
is a polyhedron with at most $2^n$ facets. This is the standard external
input in this argument. Since $E_\tau$ is open, its inclusion in that
polyhedron places it in the interior. Its facet complements give the
desired admissible cover. If $E_\tau$ is empty, the whole lattice is one
admissible class.
\end{proof}

One constrained variant also follows from the same standard input.
Suppose $\phi$ is lower semicontinuous, $\phi(z)\ge\OPT$ for every
$z\in\Z^n$, $\tau<\OPT$, and the closed sublevel set
$C_\tau=\{\phi\le\tau\}$ is bounded. It is compact, convex, and
integer-free. If nonempty, its distance $\eta$ from $\Z^n$ is positive.
The open convex set $C_\tau+B^\circ(0,\eta/2)$ is lattice-free. Facet
complements of a maximal lattice-free polyhedron enclosing this set
cover the lattice and avoid $C_\tau$, so
$\kappa_\tau(\Z^n)\le2^n$ here as well. Lower semicontinuity and the
strict threshold are needed in this compact-neighborhood argument.

There is also an elementary parity bound on midpoint cliques whenever
$\phi(z)\ge\tau$ at every integer point: two points congruent modulo two
have an integer midpoint and cannot conflict. Every midpoint clique then
has at most $2^n$ vertices. This does not bound the segment graph or the
class number without additional hypotheses.

\subsection{Random-lattice moments and scale bounds}
\label{app:lattice:cvp-proof}

The probabilistic proofs use the following standard external theorem, with
the normalization of \eqref{lattice:cvp-model}: for a nonnegative integrable
test function $h$ on $\R^n$ and Haar-random unimodular $L$, $n\ge2$,
\begin{equation}
 \E_L\sum_{a\in L\setminus\{0\}}h(a)=\int_{\R^n}h(a)\,da.
 \label{lattice:siegel}
\end{equation}
This is Siegel's mean-value theorem; see
\citet[Theorem~3.2 and Remark~3.3]{skenderi2021RandomLatticesSiegel} for
this normalization and test-function class. All moment statements used here follow
from \eqref{lattice:siegel} and the elementary torus calculation below.

Let $A\subset\R^n$ be bounded and measurable, with volume $V$, and write
$N_A=\#\{v\in L:v-t\in A\}$. A fundamental domain $D$ for $L$ has
volume one. Its translates $v-D$, $v\in L$, tile $\R^n$ up to null sets.
Tonelli's theorem gives, for any nonnegative measurable $h$,
\begin{equation}
 \E_t\sum_{v\in L}h(v-t)
 =\sum_{v\in L}\int_{v-D}h(x)\,dx=\int h(x)\,dx.
 \label{lattice:unfolding}
\end{equation}
For two such functions $h_1,h_2$, write the second lattice point as
$v+a$ and apply \eqref{lattice:unfolding} to obtain
\begin{equation}
 \E_t\left(\sum_vh_1(v-t)\right)
        \left(\sum_wh_2(w-t)\right)
 =\sum_{a\in L}\int h_1(x)h_2(x+a)\,dx.
 \label{lattice:pair-unfolding}
\end{equation}
With both functions equal to $\mathbf1_A$, the $a=0$ term is $V$.
Applying \eqref{lattice:siegel} to the other terms gives
\begin{equation}
 \E_{L,t}N_A=V,
 \qquad \E_{L,t}N_A^2=V+V^2,
 \qquad \operatorname{Var}_{L,t}(N_A)=V.
 \label{lattice:moments}
\end{equation}
The integral of the intersection-volume function equals $V^2$, again by
Tonelli. Similarly, the expected number of ordered distinct pairs in $A$
at distance less than $s$ is
\begin{equation}
 \int_{\|a\|<s}\vol(A\cap(A-a))\,da.
 \label{lattice:close-pairs}
\end{equation}

Since a ball of radius $r g_n$ has volume $r^n$, Markov's inequality and
\eqref{lattice:moments} give
\begin{align}
 \Pbb\{\dist(t,L)<(1-\delta)g_n\}&\le(1-\delta)^n,
 \label{lattice:radius-low}\\
 \Pbb\{\dist(t,L)>(1+\delta)g_n\}&\le(1+\delta)^{-n}.
 \label{lattice:radius-high}
\end{align}
For the second inequality, a ball of volume $V=(1+\delta)^n$ has no
lattice point on the event in question; Chebyshev bounds that void event
by $V/V^2$. Nonzero lattice vectors occur in opposite pairs, so
\eqref{lattice:siegel} and Markov also give
\begin{equation}
 \Pbb\{\lambda_1(L)<(1-\delta)g_n\}
 \le\tfrac12(1-\delta)^n.
 \label{lattice:short-vector}
\end{equation}
No independence between these three events is used.

\subsection{Cap occupancy and class bounds}
\label{app:lattice:cap-proof}

\begin{proof}[Proof of \cref{lattice:cap}]
Consider the finite subset of lattice points from an admissible class in
$B^\circ(t,R)$, translated by $-t$. If it is empty there is nothing to
prove. Its compact convex hull has a point $p$ of minimum norm. Since it
lies in the full class hull, admissibility gives $\|p\|\ge\sqrt\tau>0$.
For every point $u$ in the finite subset, differentiation along the
segment from $p$ to $u$ gives
$\langle p,u-p\rangle\ge0$. Put $e=p/\|p\|$. Then
$\langle e,u\rangle\ge\|p\|\ge\sqrt\tau$, and hence
\begin{equation*}
 \|u-\sqrt\tau e\|^2
 =\|u\|^2-2\sqrt\tau\langle e,u\rangle+\tau
 <R^2-\tau\le\lambda_1(L)^2/2.
\end{equation*}
Distinct such points are separated by at least $\lambda_1(L)$. Their
vectors $v,w$ from the center $\sqrt\tau e$ therefore obey
\begin{equation*}
 \langle v,w\rangle
 =\tfrac12(\|v\|^2+\|w\|^2-\|v-w\|^2)<0.
\end{equation*}
If one vector is zero, it is the only point, since the enclosing ball has
radius less than $\lambda_1(L)$. Otherwise there are at most $n+1$
vectors with pairwise strictly negative inner products. To prove this,
suppose $n+2$ existed. Their affine dependence has coefficients summing
to zero, with both positive and negative coefficients. Splitting these
coefficients gives disjoint nonempty index sets $S,T$ and positive
coefficients with
$w_0=\sum_{i\in S}\alpha_i v_i=\sum_{j\in T}\beta_jv_j$.
Taking the inner product of the two expressions yields
$\|w_0\|^2=\sum_{i\in S,j\in T}\alpha_i\beta_j
\langle v_i,v_j\rangle<0$, a contradiction. Finally apply
\eqref{lattice:count-bound} to all the lattice points in the outer ball.
\end{proof}

\begin{proof}[Lower bound in \cref{lattice:cvp}]
Take the outer ball of radius $g_nR_-$, of volume $V=R_-^n$. On the events
\begin{equation*}
 \OPT\ge(1-\delta)^2g_n^2,
 \quad\lambda_1(L)\ge(1-\delta)g_n,
 \quad N_{B^\circ(0,g_nR_-)}\ge V/2,
\end{equation*}
we have $\tau>0$ and
\begin{equation*}
 g_n^2R_-^2-\tau
 \le g_n^2\bigl[\tfrac32(1-\delta)^2-q
                     -(1-\delta)^2+q\bigr]
 =\tfrac12(1-\delta)^2g_n^2\le\lambda_1(L)^2/2.
\end{equation*}
The cap lemma gives $\kappa_\tau\ge V/[2(n+1)]$. The first two failure
probabilities total at most $\tfrac32(1-\delta)^n$ by
\eqref{lattice:radius-low} and \eqref{lattice:short-vector}. Chebyshev and
\eqref{lattice:moments} bound the third by $4/V$. Their union bound proves
\eqref{lattice:cvp-lower}.
\end{proof}

\begin{lemma}[An elementary spherical cap cover]
\label{lattice:sphere-cover}
For $\theta$ in a compact subinterval of $(0,\pi/2)$, the unit sphere
$S^{n-1}$ has a cover by at most
$C n^2\log(n+1)(\sin\theta)^{-n}$ spherical caps of angular radius
$\theta$, with $C$ uniform over that interval.
\end{lemma}

\begin{proof}
Take a maximal chordal $1/n$-separated subset $Y$ of the sphere. The
Euclidean balls of radius $1/(2n)$ about its points are disjoint and lie
in the ball of radius $1+1/(2n)$. Comparing their volumes gives
$|Y|\le(1+2n)^n$. Maximality makes it a chordal $1/n$-net. For unit
vectors, an angular separation $\alpha$ has chord length
$2\sin(\alpha/2)$, and $\arcsin x\le\pi x/2$ on $[0,1]$; the angular
mesh is therefore at most $\pi/(2n)$.

Let $\varphi=\theta-\pi/(2n)$. For sufficiently large $n$ it exceeds
$1/n$ and is less than $\pi/2$. A uniform random cap center covers a fixed
point within angle $\varphi$ with probability
\begin{equation*}
 \sigma_n(\varphi)
 =\frac{\int_0^\varphi\sin^{n-2}u\,du}
        {\int_0^\pi\sin^{n-2}u\,du}
 \ge\frac{\sin^{n-2}(\varphi-1/n)}{\pi n}.
\end{equation*}
For $M$ independent centers, a fixed net point is missed with probability
at most $e^{-M\sigma_n(\varphi)}$. Choose
$M>n\log(1+2n)/\sigma_n(\varphi)$. A union bound shows that some choice
of centers covers every net point within angle $\varphi$. The angular
mesh then covers the full sphere within $\theta$.
Uniformly over the specified angle interval,
\begin{equation*}
 \sin^{-n}\left(\theta-\frac{\pi/2+1}{n}\right)
 \le C_0\sin^{-n}\theta
\end{equation*}
by the mean-value bound for $\log\sin\theta$ on a slightly larger compact
interval. Combining the estimates gives the stated $M$ bound. Increase
$C$ to cover the finitely many smaller dimensions, using an ordinary
finite angular net there.
\end{proof}

\begin{proof}[Upper bound in \cref{lattice:cvp}]
If $\tau\le0$, the nonnegative quadratic objective makes the entire
lattice one admissible class. Suppose $\tau>0$. In Euclidean coordinates
centered at the target, each halfspace
\begin{equation*}
 H_e=\{u:\langle e,u\rangle\ge\sqrt\tau\},\qquad\|e\|=1,
\end{equation*}
is admissible. Every lattice point $u=v-t$ has norm at least
$\sqrt{\OPT}\ge\sqrt\tau$, so the halfspace with $e=u/\|u\|$ covers it.

Put $a_+=(1+\delta)^2-q>0$ and $R_+^2=1+a_+$. On
$\OPT\le(1+\delta)^2g_n^2$, $\tau\le a_+g_n^2$. Cover each point with
$\|u\|<g_nR_+$ by its own halfspace. The number of these points has mean
$R_+^n$, so it is at most $e^{\gamma n}R_+^n$ except on an event of
probability $e^{-\gamma n}$.

Choose $\theta$ by $\sin\theta=1/R_+$. A far point
$\|u\|\ge g_nR_+$ whose direction is within $\theta$ of $e$ satisfies
\begin{equation*}
 \langle e,u\rangle\ge g_nR_+\cos\theta
 =g_n\sqrt{a_+}\ge\sqrt\tau.
\end{equation*}
A cap cover supplied by \cref{lattice:sphere-cover} therefore covers all
far lattice points by at most
$C n^2\log(n+1)R_+^n$ admissible halfspaces. When
$0\le q\le1/2$ and $0<\delta<0.1$, these angles lie in a fixed compact
subinterval of $(0,\pi/2)$, so $C$ is absolute. Combining the near and
far covers gives \eqref{lattice:cvp-upper}. Equation
\eqref{lattice:radius-high} supplies the remaining failure probability.

For \eqref{lattice:cvp-exponents}, take $\delta_n\downarrow0$ and
$\gamma_n\downarrow0$ slowly enough that $n\delta_n,n\gamma_n\to\infty$,
for example $n^{-1/3}$, and use $q_n\to q_0<1/2$. The lower outer-ball
volume is exponential, and its base tends to $\sqrt{3/2-q_0}>1$.
The upper base tends to $\sqrt{2-q_0}$, while its prefactor has logarithm
$o(n)$. The explicit failure bounds tend to zero, proving the assertion.
\end{proof}

\subsection{Midpoint clique estimate}
\label{app:lattice:clique-proof}

\begin{proof}[Proof of \cref{lattice:cvp-clique}]
Divide all Euclidean lengths by $g_n$. A ball of normalized radius $r$
then has volume $r^n$ under the rescaled, unit-ball-normalized measure.
Write $R^2=\rho r_0^2$ and $s^2=4(\rho-1)r_0^2$. On
$\OPT/g_n^2\ge(1-\delta)^2$, a pair of translated lattice points $u,w$
is conflicting if $\|u+w\|<2r_0$. If $\|u\|,\|w\|<R$ and they are
not conflicting, the parallelogram identity yields
\begin{equation*}
 4r_0^2\le\|u+w\|^2
 =2\|u\|^2+2\|w\|^2-\|u-w\|^2
 <4R^2-\|u-w\|^2.
\end{equation*}
Thus every nonconflicting pair in this ball has distance less than $s$.
Let $N$ be its point count and let $M$ count all unordered distinct pairs
in the ball at distance less than $s$. Removing one endpoint from each
such pair leaves a midpoint clique of size at least $N-M$ on the stated
radius event.

Equations \eqref{lattice:close-pairs} and symmetry give
$\E M=J/2$, where, in the normalized coordinates,
\begin{equation*}
 J=\int_{\|a\|<s}\vol(B_R\cap(B_R-a))\,da
   =\int_{B_R}\vol(B_R\cap B(x,s))\,dx.
\end{equation*}
The normalized measure here means the original physical volume after
rescaling; its unit ball has measure one. For $z\in B_R\cap B(x,s)$ and
$0\le\eta\le1$, averaging the squared-distance inequalities gives
\begin{equation*}
 \|z-\eta x\|^2+\eta(1-\eta)\|x\|^2
 <(1-\eta)R^2+\eta s^2.
\end{equation*}
Consequently the intersection volume is at most
$(A-b\|x\|^2)^{n/2}$, where
\begin{equation*}
 \eta=\frac{2-\rho}{\rho},\qquad
 A=(1-\eta)R^2+\eta s^2,\qquad b=\eta(1-\eta),
 \qquad a^2=A-bR^2=\frac{4(\rho-1)}\rho r_0^2.
\end{equation*}
The function $r^{n-1}(A-br^2)^{n/2}$ is increasing for $0\le r\le R$.
Indeed, its logarithmic derivative is nonnegative exactly when
$(n-1)A\ge(2n-1)br^2$; it suffices to check $r=R$, giving
$bR^2\le(1-1/n)a^2$. Substitution reduces this to $\rho\ge2/n$,
which holds for $n\ge2$ and $\rho>1$. Radial integration therefore gives
\begin{equation*}
 J\le n\int_0^R r^{n-1}(A-br^2)^{n/2}\,dr
 \le nR^n a^n=nVa^n.
\end{equation*}
Chebyshev gives $\Pbb\{N<V/2\}\le4/V$, and Markov gives
\begin{equation*}
 \Pbb\{M\ge V/4\}\le\frac{4\E M}{V}
 \le2na^n
 \le2n\left(\frac{4(\rho-1)}\rho\right)^{n/2},
\end{equation*}
using $r_0\le1$. Together with \eqref{lattice:radius-low}, these prove the
displayed probability bound and the clique size $V/4$.

For $q_n=o(1)$, choose
$\delta_n=\max\{q_n,n^{-1/3}\}$ and
$\rho_n=4/3-n^{-1/3}$ for sufficiently large $n$. Their limiting outer
volume base is $\sqrt{4/3}$; all three failure bounds tend to zero.
Polynomial prefactors and the factor $1/4$ give only $o(n)$ in the
logarithm, proving the limiting exponent.
\end{proof}

\subsection{The two-feature curvature construction}
\label{app:lattice:gadget}

For a block with feature matrix $X=[u\ v]$, the projected perspective
value is
\begin{equation}
 g(z)=\inf_\beta\left\{\|y-X\beta\|^2+
                   \lambda\sum_{i=1}^2\frac{\beta_i^2}{z_i}\right\}
 =y^T(I+X\diag(z)X^T/\lambda)^{-1}y,
 \label{lattice:perspective-identity}
\end{equation}
where $\beta_i^2/0$ is zero for $\beta_i=0$ and $+\infty$ otherwise.
To prove the identity, restrict to the positive coordinates of $z$ and
minimize the resulting strictly convex quadratic. Its value is
$y^Ty-y^TX_S(X_S^TX_S+\lambda D^{-1})^{-1}X_S^Ty$, where
$D=\diag(z_S)$. Multiplying the matrices verifies that this equals the
right side of \eqref{lattice:perspective-identity}; zero coordinates
contribute nothing. A useful direct proof of convexity is the dual formula
\begin{equation*}
 g(z)=\sup_w\left\{2y^Tw-\|w\|^2-
                    \lambda^{-1}\sum_i z_i(x_i^Tw)^2\right\}.
\end{equation*}
The maximum is attained at
$w=(I+X\diag(z)X^T/\lambda)^{-1}y$. It is a supremum of affine functions
of $z$, so $g$ is convex. At binary $z$, the original minimization is
exact ridge regression on its selected columns. Orthogonal blocks make
these formulas add as in \eqref{lattice:block-function}.

\begin{proof}[Proof of \cref{lattice:curvature-gadget}]
Choose $\pi$ strictly between the two sides of the first inequality in
\eqref{lattice:gadget-hypotheses}. Ridge regression cannot worsen when a
feature is added, so $g_{11,j}\le g_j^*$. Hence $\pi>0$. For an integer
block selection $p\in\{0,1\}^2$,
\begin{equation}
 g_j(p)\ge g_j^*-\pi(\mathbf1^Tp-1).
 \label{lattice:block-minorant}
\end{equation}
This holds by the choice of $\pi$ for the zero- and two-feature
selections, strictly in both cases, and by the definition of $g_j^*$ for
one-feature selections. Summing gives, for every support with at most $k$
features,
\begin{equation*}
 g(z)\ge\sum_jg_j^*-\pi(\mathbf1^Tz-k)\ge\sum_jg_j^*.
\end{equation*}
Choosing the better feature in each block attains equality. Equality
requires exactly one feature per block. If every $\Delta_j>0$, each
block's better feature is unique, proving uniqueness of the optimum.

For two different one-per-block supports, let $D$ be the nonempty set of
blocks on which they differ. Their midpoint selects both features with
weight one half on $D$ and agrees with one binary feature elsewhere. It
belongs to the root relaxation and has value at most
\begin{equation*}
 \sum_{j\in D}(g_j^*-d_j)
 +\sum_{j\notin D}(g_j^*+\Delta_j)
 \le\OPT-\min_jd_j+\sum_j\Delta_j
 <\OPT-\varepsilon.
\end{equation*}
Thus all $2^k$ supports form a midpoint clique of feasible indicators.
The class-number lower bound and the binary removal accounting give the
claimed leaves and nodes.

For the stronger relaxation, every point of a block's original
mixed-integer objective epigraph satisfies the affine inequality
$t_j\ge g_j^*-\pi(\mathbf1^Tz_j-1)$, by
\eqref{lattice:block-minorant}. It remains valid on the closed convex
hull of that epigraph. Summing it and using the root budget gives
$\sum_jt_j\ge\sum_jg_j^*=\OPT$. The optimal binary support is feasible
for the hull relaxation, so its root value is exactly $\OPT$.

It remains to construct a uniform rational family. Take $\lambda=1$,
$u=(1,0)$, $v=(3/5,4/5)$, and baseline response $y=u+v$. Direct inversion
in \eqref{lattice:perspective-identity} gives
\begin{equation*}
 g_{00}=\frac{16}{5},\quad
 g_{10}=g_{01}=\frac{48}{25},\quad
 g_{11}=\frac{16}{13},\quad
 g(1/2,1/2)=\frac{16}{9}.
\end{equation*}
The midpoint deficit is $d_0=32/225>0$, and the difference between the
first and second feature improvements is $192/325>0$.
Choose one sufficiently small positive rational $\eta$, independent of
$k$, and let
\begin{equation}
 s_j=1+\eta j/k,\qquad e_j=\eta j/k^3,
 \qquad y_j=s_j(u+v)+e_ju,
 \quad 1\le j\le k.
 \label{lattice:rational-family}
\end{equation}
All these responses lie in a common neighborhood of $u+v$. Continuity
of the four binary values and the midpoint value gives, after choosing
$\eta$ small enough, the global first inequality in
\eqref{lattice:gadget-hypotheses} and $d_j\ge d_0/2$ for every $j,k$.
This is a global inequality across blocks: shrink the neighborhood until
its smallest zero-to-one improvement exceeds its largest one-to-two
improvement, using the positive baseline margin.

Sherman--Morrison gives
$g_{10,j}=\|y_j\|^2-(u^Ty_j)^2/2$ and
$g_{01,j}=\|y_j\|^2-(v^Ty_j)^2/2$. Their difference is strictly positive
in favor of $u$, because
$u^Ty_j-v^Ty_j=(1-3/5)e_j>0$ and both inner products are positive.
On the same compact neighborhood, their sum is bounded, so
$0<\Delta_j\le C_0e_j$ for a constant $C_0$ independent of $k,j$.
Consequently
\begin{equation*}
 \sum_j\Delta_j\le C_0\eta\frac{k+1}{2k^2}\le C_0\eta.
\end{equation*}
Choose $\eta$ still smaller so that $C_0\eta<d_0/4$. Every tolerance
$0\le\varepsilon<d_0/4$ now satisfies the second inequality in
\eqref{lattice:gadget-hypotheses}, and every $\Delta_j$ is positive.
The feature coordinates and responses are rational with denominators
polynomial in $k$; their bit lengths and the total instance encoding
length are polynomial in $k$. This proves the final assertion.
\end{proof}
